\documentclass{book}
\usepackage{amsmath}
\usepackage{bm}
\usepackage{colortbl}
\usepackage{graphicx}
\usepackage{caption}
\usepackage{fullpage}
\usepackage{afterpage}
\usepackage{multirow}
\usepackage[nodisplayskipstretch]{setspace}
\usepackage{booktabs}
\usepackage{gensymb}
\usepackage{parskip}
\usepackage{xr}
\usepackage{pdflscape}
\usepackage[inline]{enumitem}
\usepackage{wrapfig}
\usepackage{longtable}
\usepackage{subfig}
\usepackage[natbib=true, style=numeric-comp,subentry,backend=biber,sorting=none]{biblatex}
\addbibresource{citation/Thesis.bib}
\usepackage[hidelinks]{hyperref}
\usepackage{geometry}
\geometry{
	top=1in,            % <-- you want to adjust this
	inner=1in,
	outer=1in,
	bottom=1in,
	headheight=3ex,       % <-- and this
	headsep=2ex,          % <-- and this
}
\usepackage{fancyhdr}
\usepackage{fixltx2e}
\usepackage[acronym,nomain,toc,makeindex ]{glossaries}
\usepackage{cleveref} %This need to be the last package for it to work properly 


\pagestyle{fancy}

\title{Genetic and Environmental risk factors of Schizophrenia and Autism}
\date{\today}
\author{\href{mailto:choishingwan@gmail.com}{Choi Shing Wan}\\
	\includegraphics[width=0.5\textwidth]{hkuLogo.jpg}}
\singlespacing
\renewcommand*\contentsname{Contents}
\onehalfspacing
%\doublespacing

%\includeonly{environmental_risk/er_chapter,supplementary_materials}

\makeglossary
\newacronym{SCZ}{SCZ}{Schizophrenia}
\newacronym{ASD}{ASD}{Autism Spectrum Disorder}
\newacronym[longplural={Genome Wide Association Studies}]{GWAS}{GWAS}{Genome Wide Association Study}
\newacronym{SNP}{SNP}{Single Nucleotide Polymorphism}
\newacronym{LD}{LD}{Linkage Disequilibrium}
\newacronym{PGS}{PGS}{Polygenic Risk Score}
\newacronym{tSVD}{tSVD}{Truncated Singular Value Decomposition}
\newacronym{SVD}{SVD}{Singular Value Decomposition}



\raggedbottom %Remove it before printing as this is something to do with global settings. Can make each page look uneven but more dense. 
\makeindex
\begin{document}\thispagestyle{empty}
\pagestyle{empty}

%\maketitle
\begin{titlepage}
	\begin{center}
		\vspace*{1cm}
		
		\Huge
		\textbf{Heritability Estimation and Risk Prediction in Schizophrenia}
		
		\vspace{0.5cm}
		\LARGE
		
		\vspace{1.5cm}
		
		\textbf{\href{mailto:choishingwan@gmail.com}{Choi Shing Wan}}
		
		\vfill
		
		A thesis submitted in partial fulfillment of the requirements for \\
		the Degree of Doctor of Philosophy
		
		\vspace{0.8cm}
		
		\includegraphics[width=0.4\textwidth]{hkuLogo.jpg}
		
		\Large
		Department of Psychiatry\\
		University of Hong Kong\\
		Hong Kong\\
		\today
		
	\end{center}
\end{titlepage}


\frontmatter 

	\cleardoublepage
	\phantomsection
	\addcontentsline{toc}{chapter}{Declaration}
	\chapter*{Declaration}
	\cleardoublepage
	\phantomsection
	\addcontentsline{toc}{chapter}{Acknowledgments}
	\chapter*{Acknowledgements}
	\cleardoublepage
	\phantomsection
	\printglossary[title=Abbreviations,toctitle=Abbreviations]
	\cleardoublepage
	\phantomsection
	%To generate the correct abbreviations, use alt+shift+F1 twice before using F1
	
	\cleardoublepage
	\phantomsection
	\addcontentsline{toc}{chapter}{Contents}
	\begin{singlespace}
		\tableofcontents
	\end{singlespace}
\mainmatter
	\chapter*{Introduction}
		\addcontentsline{toc}{chapter}{Introduction}
	\pagestyle{fancy}
	\chapter*{Some considerations}
	\begin{enumerate}
		\item PRSice requires the phenotype to aid its selection (More information= stronger)
		\item It seems like LDSC doesn't necessary perform badly in oligogenic situation.
		Rather, it is that when the trait is oligogenic, it is more likely for LDSC to behaviour in a strange way.
		\item For each condition: extreme phenotype, quantitative trait, case control, we can have a separated review. 
		Discuss on the benefits and challenges of each condition and the method we deal with them.
		So we can have two chapters (case control, quantitative trait) where extreme phenotype can be a big subsection within quantitative trait.
		\item For each chapter, there will be this introduction (review on the method), our methodology (Calculation, implementation and also simulation), result (the simulation result). 
		Then we can have the application (PGC, network)
	\end{enumerate}
	
	\chapter{Literature Review}
	\section{Twin Studies}
	Should briefly talk about how Twin modeling was used for finding the GE contribution.
	Should also mention the ACE model.
	At the end, we can talk about the heritability estimates of SCZ and AD
	\section{Searching for Genetic Variants}
	\subsection{Role of Common Variants}
	\subsubsection{Genome Wide Association Study}
	Should talk about what is GWAS and how it is used.
	Should also talk about the current GWAS studies in SCZ and AD
	\subsection{Role of Rare Variants}
	\subsubsection{Exome Sequencing}
	Similar to the GWAS.
	Talk about the Pros and Cons.
	Need to briefly mention the Denovo paper and Shaun's paper.					
	\subsubsection{Whole Genome Sequencing}
	Very very brief description of WGS and the current status.
	\section{Narrow Sense Heritability}
	\section{Risk Prediction}
	\section{Summary}
	
	\chapter{Heritability Estimation}
	This chapter should be used in similar way as the general method section in Clara's thesis. Considering that the subsequent chapters all rely on this implementation.
	\section{Introduction}
	\section{Methodology}
		The work in this chapter were done in collaboration with my colleagues who have kindly provide their support and knowledges to make this piece of work possible.
		
		\subsection{Heritability Estimation}
			The narrow-sense heritability is defined as 
			$$
				h^2 = \frac{\mathrm{Var}(X)}{\mathrm{Var}(Y)}
			$$
			where $\mathrm{Var}(X)$ is the variance of the genotype and $\mathrm{Var}r(Y)$ is the variance of the phenotype.
			In a \gls{GWAS}, regression were performed between the \glspl{SNP} and the phenotypes, giving
			\begin{equation}
				Y=\beta X+\epsilon
				\label{eq:standardRegress}
			\end{equation}
			where $Y$ and $X$ are the standardized phenotype and genotype respectively. 
			$\epsilon$ is then the error term, accounting for the non-genetic elements contributing to the phenotype (e.g. Environment factors).
			Based on \cref{eq:standardRegress}, one can then have
			\begin{align}
				\mathrm{Var}(Y) = \mathrm{Var}(\beta X)+ \mathrm{Var}(\epsilon) \nonumber\\
				\mathrm{Var}(Y) = \beta^\mathrm{Var}(X) \nonumber\\
				\beta^2\frac{\mathrm{Var}(X)}{\mathrm{Var}(Y)}= 1
				\label{eq:betaHeri}
			\end{align}
			$\beta^2$ is then considered as the portion of phenotype variance explained by the variance of genotype, which can also be considered as the narrow-sense heritability of the phenotype.
					
			A challenge in calculating the heritability from \gls{GWAS} data is that usually only the test-statistic or p-value were provided and one will not be able to directly calculate the heritability based on \cref{eq:betaHeri}. In order to estimation the heritability of a trait from the \gls{GWAS} test-statistic, we first observed that when both $X$ and $Y$ are standardized, $\beta^2$ will be equal to the coefficient of determination ($r^2$). Then, based on properties of the Pearson product-moment correlation coefficient:
			\begin{equation}
				r = \frac{t}{\sqrt{n-2+t^2}}
				\label{eq:pearsonProduct}
			\end{equation}
			where $t$ follows the student-t distribution and $n$ is the number of samples.
			One can then obtain the $r^2$ by taking the square of \cref{eq:pearsonProduct}
			\begin{equation}
				r^2 = \frac{t^2}{n-2+t^2}
				\label{eq:oriRSquared}
			\end{equation}
			It is observed that $t^2$ will follow the F-distribution and when $n$ is big, $t^2$ will converge into $\chi^2$ distribution.
			
			When the effect size is small and $n$ is big, $r^2$ will be approximately $\chi^2$ distributed with mean $\sim 1$. 
			We can then approximate \cref{eq:oriRSquared} as
			\begin{equation}
				r^2= \frac{\chi^2}{n}
				\label{eq:approxChi}
			\end{equation}
			and define the \emph{observed} effect size of each \gls{SNP} to be
			\begin{equation}
			f=\frac{\chi^2-1}{n}
			\label{eq:observedEffect}
			\end{equation}
			
			When there are \gls{LD} between each individual \glspl{SNP}, the situation will become more complicated as each \glspl{SNP}' observed effect will contains effect coming from other \glspl{SNP} in \gls{LD} with it. 
			\begin{equation}
			f_{observed} = f_{true}+f_{LD}
			\label{eq:conceptF}
			\end{equation}
			
			To account for the \gls{LD} structure, we first assume our phenotype $\boldsymbol{Y}$ and genotype $\boldsymbol{X}=(X_1,X_2,\dots,X_m)^t$ are standardized and that
			\begin{align*}
				\boldsymbol{Y}\sim f(0,1) \\
				\boldsymbol{X}\sim f(0,\boldsymbol{R})
			\end{align*}
			Where $\boldsymbol{R}$ is the \gls{LD} matrix between \glspl{SNP}.
			
			We can then express \cref{eq:standardRegress} in matrix form:
			\begin{align}
				\boldsymbol{Y}=\boldsymbol{\beta}^t\boldsymbol{X}+\epsilon
				\label{eq:matrixRegress}
			\end{align}
			Definition of heritability will then become
			\begin{align}
				Heritability& = \frac{\mathrm{Var}(\boldsymbol{\beta}^t\boldsymbol{X})}{\mathrm{Var}(\boldsymbol{Y})} \nonumber\\
				&=\mathrm{Var}(\boldsymbol{\beta}^t\boldsymbol{X})
			\end{align}
			If we then assume now that $\boldsymbol{\beta} = (\beta_1, \beta_2,\dots,\beta_m)^t$ has distribution
			\begin{align*}
				\boldsymbol{\beta}&\sim f(0,\boldmath{H})\\
				\boldsymbol{H}&=diag(\boldsymbol{h})\\
				\boldsymbol{h}&=(h_1^2,h_2^2,\dots,h_m^2)^t
			\end{align*}
			where $\boldsymbol{H}$ is the variance of the true effect. 
			It is shown that heritability can be expressed as %The later part was gone because that will contains E(\beta) which = 0
			\begin{align}
			\mathrm{Var}(\boldsymbol{\beta}^t\boldsymbol{X}) &= \mathrm{E}_X\mathrm{Var}_{\beta|X}(\boldsymbol{X}^t\boldsymbol{\beta})+\mathrm{Var}_X\mathrm{E}_{(\beta|X)}(\boldsymbol{\beta}^2\boldsymbol{X}) \nonumber\\
			&=\mathrm{E}_X(\boldsymbol{X}^t\boldsymbol{\beta\beta}^T\boldsymbol{X}) \nonumber\\ 
			&= \mathrm{E}_X(\boldsymbol{X}^t\boldsymbol{HX}) \nonumber\\
			&= \mathrm{E}(\boldsymbol{X})^t\boldsymbol{H}\mathrm{E}(\boldsymbol{X})+\mathrm{Tr}(\mathrm{Var}(\boldsymbol{X}\boldsymbol{H})) \nonumber\\
			&=\mathrm{Tr}(\mathrm{Var}(\boldsymbol{X}\boldsymbol{H})) \nonumber\\
			&=\sum_ih_i^2
			\label{eq:proveHerit}
			\end{align}
			
			Now if we consider the covariance between \gls{SNP} i ($X_i$) and $Y$, we have
			\begin{align}
			 \mathrm{Cov}(\boldsymbol{X}_i,\boldsymbol{Y}) &= \mathrm{Cov}(\boldsymbol{X}_i,\boldsymbol{\beta}^t\boldsymbol{X}+\epsilon) \nonumber\\
			 &=\mathrm{Cov}(\boldsymbol{X}_i,\boldsymbol{\beta}^t\boldsymbol{X}) \nonumber\\
			 &=\sum_j{\mathrm{Cov}(\boldsymbol{X}_i,\boldsymbol{X}_j)\boldsymbol{\beta}_j} \nonumber\\
			 &=\boldsymbol{R}_i\boldsymbol{\beta}_j
			 \label{eq:covPhenoTrue}
			\end{align}
			As both $X$ and $Y$ are standardized, the covariance will equal to the correlation and we can define the correlation between \gls{SNP} i and $Y$ as
			\begin{equation}
				\rho_i = \boldsymbol{R}_i\boldsymbol{\beta}_j
				\label{eq:corPhenoTrue}
			\end{equation}
			In reality, the \emph{observed} correlation usually contains error. 
			Therefore we define the \emph{observed} correlation to be
			\begin{equation}
			\hat{\rho_i} = \rho_i+\frac{\epsilon_i}{\sqrt{n}}
			\label{eq:obsPheno}
			\end{equation}
			for some error $\epsilon_i$. 
			The distribution of the correlation coefficient about the true correlation $\rho$ is approximately
			$$
				\hat{\rho_i}\sim f(\rho_i, \frac{(1-\rho^2)^2}{n})
			$$
			By making the assumption that $\rho_i$ is close to 0 for all $i$, we have 
			\begin{align*}
				\mathrm{E}(\epsilon_i|\rho_i)&\sim 0\\
				\mathrm{Var}(\epsilon_i|\rho_i)&\sim 1
			\end{align*}
			We then define our $z$-statistic and $\chi^2$-statistic as
			\begin{align*}
				z_i &= \hat{\rho_i}\sqrt{n} \\
				\chi^2 &= z_i^2\\
				&=\hat{\rho_i}^2n
			\end{align*}
			From \cref{eq:obsPheno} and \cref{eq:corPhenoTrue}, $\chi^2$ can then be expressed as
			\begin{align*}
			\chi^2&=\hat{\rho}^2n\\
			&=n(\boldsymbol{R}_i\boldsymbol{\beta}_j+\frac{\epsilon_i}{\sqrt{n}})^2
			\end{align*}
			The expectation of $\chi^2$ is then
			\begin{align*}
			\mathrm{E}(\chi^2) &= n(\boldsymbol{R}_i\boldsymbol{\beta\beta}^t\boldsymbol{R}_i+2\boldsymbol{R}_i\boldsymbol{\beta}\frac{\epsilon_i}{\sqrt{n}}+\frac{\epsilon_i^2}{n}) \\
			&= n\boldsymbol{R}_i\boldsymbol{H}\boldsymbol{R}_i+1
			\end{align*}
			To derive least square estimates of $h_i^2$, we need to find $\hat{h_i^2}$ which minimizes
			\begin{align*}
				\sum_i(\chi_i^2-\mathrm{E}(\chi_i^2))^2&=\sum_i(\chi_i^2-(n\boldsymbol{R}_i\boldsymbol{H}\boldsymbol{R}_i+1))^2 \\
				&=\sum_i(\chi_i^2-1-n\boldsymbol{R}_i\boldsymbol{H}\boldsymbol{R}_i)^2 
			\end{align*}
			If we define 
			\begin{equation}
			f_i= \frac{\chi_i^2-1}{n}
			\label{eq:defineF}
			\end{equation}
			we got
			\begin{align}
			\sum_i(\chi_i^2-\mathrm{E}(\chi_i^2))^2&=\sum_i(f_i-\boldsymbol{R}_i\boldsymbol{H}\boldsymbol{R}_i)^2 \nonumber\\
			&=\boldsymbol{ff}^t-2\boldsymbol{f}^t\boldsymbol{R_{sq}\hat{h}}+\boldsymbol{\hat{h}}^t\boldsymbol{R_{sq}}^t\boldsymbol{R_{sq}\hat{h}}
			\label{eq:leastSquareH}
			\end{align}
			where $\boldsymbol{R_{sq}} = \boldsymbol{R}\circ\boldsymbol{R}$.
			By differentiating \cref{eq:leastSquareH} w.r.t $\hat{h}$ and set to 0, we get
			\begin{align}
				2\boldsymbol{R_{sq}}^t\boldsymbol{R_{sq}}\boldsymbol{\hat{h^2}}-2\boldsymbol{R_{sq}f}&=0 \nonumber\\
				\boldsymbol{R_{sq}}\boldsymbol{\hat{h^2}} &=\boldsymbol{f}
				\label{eq:shrekEq}
			\end{align}
			And the heritability is then defined as 
			\begin{equation}
			\hat{Heritability} = \boldsymbol{1}^t\boldsymbol{R_{sq}}^{-1}\boldsymbol{f}
			\label{eq:fullShrek}
			\end{equation}
		\subsection{Inverse of the \glsentrylong {LD} matrix}
			In order to obtain the heritability estimation, we will require to solve \cref{eq:fullShrek}. 
			If $\boldsymbol{R_{sq}}$ is of full rank and positive semi-definite, it will be straight-forward to solve the matrix equation.
			However, more often than not, the \gls{LD} matrix are rank-deficient and suffer from multicollinearity, making it ill-conditioned, therefore highly sensitive to changes or errors in the input.
			To be exact, we can view \cref{eq:fullShrek} as calculating the sum of $\boldsymbol{\hat{h^2}}$ from  \cref{eq:shrekEq}.
			This will involve solving for
			\begin{equation}
			\boldsymbol{\hat{h^2}} = \boldsymbol{R_{sq}}^{-1}\boldsymbol{f}
			\label{eq:shrekInverse}
			\end{equation}
			where an inverse of $\boldsymbol{R_{sq}}$ is observed. 
			
			In normal circumstances (e.g. when $\boldsymbol{R_{sq}}$ is full rank and positive semi-definite), one can easily solve \cref{eq:shrekInverse} using the QR decomposition or LU decomposition.
			However, when $\boldsymbol{R_{sq}}$ is ill-conditioned, the traditional decomposition method will fail.
			Even if the decomposition is successfully performed, the result tends to be a meaningless approximation to the true $\boldsymbol{\hat{h^2}}$. 
			
			Therefore, to obtain a meaningful solution, regularization techniques such as the Tikhonov Regularization (also known as Ridge Regression) and \gls{tSVD} has to be performed\cite{Neumaier1998}. 
			Arguably, there are a large variety of regularization techniques, yet the discussion of which is beyond the scope of this study. 
			In this study, we will focus on the use of \gls{tSVD} in the regularization of the \gls{LD} matrix.
			This is because the \gls{SVD} routine has been implemented in the EIGEN C++ library \cite{eigenweb}, allowing us to implement the \gls{tSVD} method without much concern with regard to the algorithm. 
			
			To understand the regularization method, we first assume
			
			
			In this study, \gls{tSVD} were used as the \gls{SVD} was implemented in the EIGEN C++ library\cite{eigenweb}, allow for a quick and simple implementation of the method.
			%Start explaining what is SVD and how to perform the tSVD. This will be complicated and slightly too difficult for me ><
			
			
			
	\subsection{Quantitative Trait}
	\subsection{Case Control Studies}
	\subsection{Extreme Phenotype Selections}
	\section{Simulation}
	\subsection{Quantitative Trait}
	\subsection{Case Control Studies}
	\subsection{Exreme Phenotype Selections}
	\section{Result}
	\section{Discussion}
	
	\chapter{Heritability of Schizophrenia}
	\section{Introduction}
	\section{Heritability Estimation}
	This will be a very simple section, focused on how to perform the heritability estimation on \acrfull{SCZ}.
	Should also tokenize the heritability into subcategories (e.g. immune, neuron, etc)
	\subsection{Methodology}
	\subsection{Result}
	\section{Brain development and Schizophrenia}
	Here we will perform the WGCNA and brain development network.
	Seeing how the whether if any brain development network were enriched with SNPs that explain the variance of phenotype
	\subsection{Methodology}
	\subsection{Result}
	\section{Discussion}
	\chapter{Heritability of Response to antipsychotic treatment}
	\section{Introduction}
	Here we try to use Beatrice's data and estimate the heritability explained in drug response.
	Should also repeat the region-wise heritability
	\section{Methodology}
	\section{Result}
	\section{Discussion}
	\chapter{Risk Prediction}
	\section{Methodology}
	We can define the traditional \gls{PGS} as
	\begin{equation}
	\hat{Y} = diag(\beta)X
	\end{equation}
	where $X$ is the standardized genotype, $\beta$ is the test-statistic calculated from other studies. 
	
	
	\subsection{Simulation}
	\section{Result}
	\section{Discussion}
	\chapter{Conclusion}
	\backmatter
	\printbibliography
	\chapter*{Appendix}
\end{document}


